\section{Formulación del problema}
\label{sec:formulacion}

La pregunta general del trabajo es la siguiente: ¿cómo estimar la tasa de desempleo en los municipios de los departamentos del Cauca y Valle del Cauca que no tienen información directa, mediante técnicas de estimación de áreas pequeñas basadas en modelos mixtos con componentes espaciales, de modo que se superen las limitaciones de cobertura y costo de los métodos tradicionales?

Las preguntas de sistematización, una por cada objetivo específico, son:

\begin{itemize}
    \item ¿Qué modelos de estimación de áreas pequeñas basados en modelos mixtos con componentes espaciales ofrecen mejor desempeño para estimar la tasa de desempleo en los municipios de los departamentos del Cauca y Valle del Cauca, con las fuentes de información disponibles (encuestas de hogares, registros administrativos, censos y otras)?

    \item ¿Cuáles variables explicativas, según los resultados de los modelos de estimación de áreas pequeñas y la evidencia previa disponible, tienen mayor incidencia en la tasa de desempleo en los municipios de los departamentos del Cauca y Valle del Cauca?

    \item ¿Cómo se explican la distribución y los patrones espaciales del desempleo en los municipios de los departamentos del Cauca y Valle del Cauca a partir de las estimaciones obtenidas mediante los modelos de estimación de áreas pequeñas?
\end{itemize}
